% !TEX root = main.tex
\section{IGA Continuum Mechanics, Discretization, and Pullback Mapping}
\label{sec:iga_discretization}

This section details the continuum elasticity formulas, strain-displacement operators, stiffness/mass element matrix assembly, Gauss quadrature transformations, and the diffeomorphic reference pullback mapping $\hat{\Omega}$.

\subsection{Linear Elasticity Constitutive Matrix $\mathbf{D}$ (Hooke's Law)}

\subsubsection{Formulas}
\paragraph{2D Plane Stress (Thin Plates / Cantilever Beams in $x-y$ Plane)}
In 2D plane stress ($\sigma_z = 0, \tau_{xz} = 0, \tau_{yz} = 0$), the stress-strain relation $\boldsymbol{\sigma} = \mathbf{D} \boldsymbol{\varepsilon}$ is governed by:
\begin{equation}
\mathbf{D}_{\text{plane stress}} = \frac{E}{1 - \nu^2} \begin{bmatrix} 
1 & \nu & 0 \\ 
\nu & 1 & 0 \\ 
0 & 0 & \frac{1 - \nu}{2} 
\end{bmatrix}
\end{equation}

\paragraph{2D Plane Strain (Thick Structures in $z$-direction, $\varepsilon_z = 0$)}
\begin{equation}
\mathbf{D}_{\text{plane strain}} = \frac{E}{(1 + \nu)(1 - 2\nu)} \begin{bmatrix} 
1 - \nu & \nu & 0 \\ 
\nu & 1 - \nu & 0 \\ 
0 & 0 & \frac{1 - 2\nu}{2} 
\end{bmatrix}
\end{equation}

\subsubsection{Mathematical Explanation \& Structural Rationale}
\begin{itemize}
    \item \textbf{Physical Quantities}: $E$ is Young's Modulus (material stiffness in MPa/Pa), and $\nu$ is Poisson's ratio (lateral contraction coupling).
    \item \textbf{Matrix Symmetry}: $\mathbf{D} = \mathbf{D}^T$ ensures energy conservation (existence of a strain energy density function $W = \frac{1}{2} \boldsymbol{\varepsilon}^T \mathbf{D} \boldsymbol{\varepsilon}$).
    \item \textbf{Shear Modulus Connection}: The bottom-right entry $\frac{E (1-\nu)}{2(1-\nu^2)} = \frac{E}{2(1+\nu)} = G$, which is the shear modulus linking shear stress $\tau_{xy}$ to shear strain $\gamma_{xy}$.
\end{itemize}

\subsubsection{Variable Dependencies}
\begin{enumerate}
    \item \textbf{$E$ (Young's Modulus)}: Scales all stress components linearly. Higher $E$ increases deformation resistance.
    \item \textbf{$\nu$ (Poisson's Ratio)}: Controls lateral expansion/contraction coupling. As $\nu \to 0.5$ (incompressible materials like rubber), plane strain denominator $(1 - 2\nu) \to 0$, causing volumetric locking.
\end{enumerate}

\subsubsection{App \& Code Alignment}
In \texttt{app/shared/fom-solver.js} and \texttt{visualizer/mesh\_convergence.py}:
\begin{itemize}
    \item Implements plane stress factor \texttt{E / (1 - nu * nu)} and matrix entries.
\end{itemize}

\hr

\subsection{Strain-Displacement Matrix $\mathbf{B}_a$}

\subsubsection{Formula}
For a 2D elastic domain, strain vector $\boldsymbol{\varepsilon} = [\varepsilon_{xx}, \varepsilon_{yy}, \gamma_{xy}]^T$ at any point is related to control node displacements $\mathbf{d} = [u_1, v_1, u_2, v_2, \dots]^T$ via:
\begin{equation}
\boldsymbol{\varepsilon}(u, v) = \sum_{a=1}^{(p+1)(q+1)} \mathbf{B}_a(u, v) \mathbf{d}_a
\end{equation}
where the $3 \times 2$ block matrix $\mathbf{B}_a$ for basis function $R_a$ is:
\begin{equation}
\mathbf{B}_a(u, v) = \begin{bmatrix} 
\dfrac{\partial R_a}{\partial x} & 0 \\[8pt] 
0 & \dfrac{\partial R_a}{\partial y} \\[8pt] 
\dfrac{\partial R_a}{\partial y} & \dfrac{\partial R_a}{\partial x} 
\end{bmatrix}
\end{equation}

\subsubsection{Mathematical Explanation \& Structural Rationale}
Derived directly from linear infinitesimal strain kinematics:
\begin{equation}
\varepsilon_{xx} = \frac{\partial u_x}{\partial x}, \quad 
\varepsilon_{yy} = \frac{\partial u_y}{\partial y}, \quad 
\gamma_{xy} = \frac{\partial u_x}{\partial y} + \frac{\partial u_y}{\partial x}
\end{equation}

\hr

\subsection{Element Stiffness Matrix $\mathbf{K}_e$}

\subsubsection{Formula}
The element stiffness matrix $\mathbf{K}_e \in \mathbb{R}^{2N_e \times 2N_e}$ ($N_e = (p+1)(q+1)$) for active element $e$ is:
\begin{equation}
\mathbf{K}_e = t_h \iint_{\Omega_e} \mathbf{B}^T \mathbf{D} \mathbf{B} \, dx \, dy = t_h \int_{u_i}^{u_{i+1}} \int_{v_j}^{v_{j+1}} \mathbf{B}^T \mathbf{D} \mathbf{B} \, |\det \mathbf{J}| \, du \, dv
\end{equation}

Using Gauss-Legendre quadrature mapping parametric knot span $[u_i, u_{i+1}] \times [v_j, v_{j+1}]$ to parent domain $[-1, 1] \times [-1, 1]$:
\begin{equation}
\mathbf{K}_e = t_h \sum_{g_u=1}^{N_g} \sum_{g_v=1}^{N_g} \mathbf{B}(\bar{\xi}_{g_u}, \bar{\eta}_{g_v})^T \mathbf{D} \mathbf{B}(\bar{\xi}_{g_u}, \bar{\eta}_{g_v}) \, |\det \mathbf{J}| \, \frac{\Delta u_i}{2} \frac{\Delta v_j}{2} w_{g_u} w_{g_v}
\end{equation}

\hr

\subsection{Element Mass Matrix $\mathbf{M}_e$}

\subsubsection{Formula}
The consistent element mass matrix $\mathbf{M}_e$ is:
\begin{equation}
\mathbf{M}_e = \rho t_h \iint_{\Omega_e} \mathbf{R}^T \mathbf{R} \, dx \, dy = \rho t_h \int_{u_i}^{u_{i+1}} \int_{v_j}^{v_{j+1}} \mathbf{R}^T \mathbf{R} \, |\det \mathbf{J}| \, du \, dv
\end{equation}
where $\mathbf{R}_a = \begin{bmatrix} R_a & 0 \\ 0 & R_a \end{bmatrix}$ and $\rho$ is mass density.

\hr

\subsection{Parameter-Dependent Pullback Mapping $\boldsymbol{\Psi}(\boldsymbol{\xi}; \boldsymbol{\mu})$}

\subsubsection{Structural Problem: Vector Space Incompatibility}
When structural parameters $\boldsymbol{\mu} = (r, x_c)^T$ change (notch radius or notch position), the physical domain $\Omega(\boldsymbol{\mu})$ changes. Collecting physical displacement snapshots $\mathbf{u}(\mathbf{x}; \boldsymbol{\mu})$ across different geometries evaluates them at different spatial locations, invalidating standard SVD/POD.

\subsubsection{Formula (Diffeomorphic Mapping to Static Reference Domain $\hat{\Omega}$)}
We introduce a smooth mapping $\boldsymbol{\Psi}: \hat{\Omega} \to \Omega(\boldsymbol{\mu})$ such that:
\begin{equation}
\mathbf{x}(\boldsymbol{\xi}; \boldsymbol{\mu}) = \hat{\mathbf{x}}(\boldsymbol{\xi}) + \mathbf{v}_{\text{map}}(\boldsymbol{\xi}; \boldsymbol{\mu})
\end{equation}
where $\hat{\Omega}$ is a fixed reference domain ($\boldsymbol{\mu} = \boldsymbol{\mu}_0$), and $\mathbf{v}_{\text{map}}(\boldsymbol{\xi}; \boldsymbol{\mu}) = \sum R_a(\boldsymbol{\xi}) \Delta \mathbf{P}_a(\boldsymbol{\mu})$ is the smooth control point displacement mapping field.

\subsubsection{Shifted Coordinate Jacobian Formula}
The total mapping Jacobian $\mathbf{J}(\boldsymbol{\xi}; \boldsymbol{\mu})$ becomes:
\begin{equation}
\mathbf{J}(\boldsymbol{\xi}; \boldsymbol{\mu}) = \mathbf{J}_0(\boldsymbol{\xi}) + \nabla_{\boldsymbol{\xi}} \mathbf{v}_{\text{map}}(\boldsymbol{\xi}; \boldsymbol{\mu})
\end{equation}

\subsubsection{Key Theoretical Consequence}
Because reference basis functions $\hat{R}_a(\boldsymbol{\xi})$ and reference mesh topology on $\hat{\Omega}$ are \textbf{static}:
\begin{itemize}
    \item All snapshot vectors $\mathbf{u}(\boldsymbol{\xi}; \boldsymbol{\mu}) \in \mathcal{V}_0(\hat{\Omega})$ reside in the \textbf{SAME unified linear vector space}!
    \item Singular Value Decomposition (SVD) and POD subspace extraction become mathematically rigorous and exact!
\end{itemize}
